\section{Training: optimizers, activations, regularization}

\begin{summary}{The rapid-fire block (Hochreiter asks these fast)}
\textsf{Optimizers:} GD, SGD, momentum, AdaGrad, RMSprop, Adam. \
\textsf{Activations:} sigmoid, tanh, ReLU, Leaky ReLU, ELU, SELU. \
\textsf{From Linz:} ELU (Clevert, Unterthiner, Hochreiter) and SELU (Klambauer et al.\ 2017). \
\textsf{Regularization:} weight decay, dropout, early stopping, data augmentation, noise, Flat
Minimum Search (Hochreiter). \
\textsf{Output layer:} linear (regression), sigmoid (binary), softmax (multi-class). \
\textsf{Why mini-batches:} good gradient estimate, fast on GPU, useful noise.
\[
\mathbf{w}\leftarrow\mathbf{w}-\eta\nabla\Remp,\qquad
\text{Adam: }\mathbf{w}\leftarrow\mathbf{w}-\eta\,\frac{\hat{\mathbf{m}}}{\sqrt{\hat{\mathbf{v}}}+\epsilon},\qquad
\text{SELU: }\lambda\begin{cases}x & x>0\\ \alpha(e^x-1)& x\le0\end{cases}
\]
with $\lambda\approx1.0507$, $\alpha\approx1.6733$.
\end{summary}

\pic{Walking down a foggy mountain. Gradient descent: look at the slope under your feet and step
downhill. Momentum: a heavy ball that keeps rolling through small bumps. Adam: shoes that adjust
their step to how steep and how noisy each direction is. Mini-batch: asking a few people for the
way instead of everybody (too slow) or one person (too noisy).}

\board{GD $\mathbf{w}\leftarrow\mathbf{w}-\eta\nabla\Remp$ \textbullet{}
momentum $\mathbf{v}\leftarrow\mu\mathbf{v}-\eta\nabla\Remp$, $\mathbf{w}\leftarrow\mathbf{w}+\mathbf{v}$ \textbullet{}
weight decay $\Remp+\frac\lambda2\lVert\mathbf{w}\rVert^2$ \textbullet{}
dropout: keep unit with prob.\ $1-p$ \textbullet{}
batch norm $\hat s=(s-\mu_B)/\sqrt{\sigma_B^2+\epsilon}$, then $\gamma\hat s+\beta$.}

\begin{qabox}[Hochreiter]{Which activation function comes from Linz? What does it do?}
\ans SELU, the scaled exponential linear unit, from Klambauer, Unterthiner, Mayr and Hochreiter
(2017); and before it ELU. SELU makes networks \emph{self-normalizing}: mean and variance of the
activations are pulled to $(0,1)$ layer by layer, without batch normalization.
\formula The map of (mean, variance) from one layer to the next has the fixed point $(0,1)$; its
Jacobian there has spectral norm $0.7877<1$, so the fixed point is stable (a contraction).
\pic{A thermostat: whatever the temperature, the layer pushes it back to the set point.}
\src{DLNN script WS24, Sec.~10.4, PDF p.~137--141.}
\end{qabox}

\begin{qabox}[Hochreiter]{Which optimizers do you know? Are there methods other than gradient descent?}
\ans GD, SGD, momentum, AdaGrad, RMSprop, Adam. Other methods: closed form (linear models),
IRLS, RProp and Quickprop, second-order methods (Newton, Gauss--Newton, Levenberg--Marquardt,
conjugate gradients, natural gradient), and gradient-free search (evolutionary algorithms).
\src{DLNN script WS24, Ch.~7, PDF p.~85--101; SS25 script ``Higher Order Optimization''.}
\end{qabox}

\begin{qabox}[Hochreiter]{Which regularization techniques do you know?}
\ans Weight decay, dropout, early stopping, data augmentation, input and weight noise, weight
sharing, multi-task learning, pruning (Optimal Brain Damage and Surgeon), and Flat Minimum Search,
which looks for wide minima that generalize better.
\formula Weight decay gradient: $\nabla\Remp+\lambda\mathbf{w}$ (the weights shrink each step).
\follow ``Why does dropout work?'' It trains an ensemble of thinned networks and stops units from
co-adapting.
\src{DLNN script WS24, Ch.~8, PDF p.~103--119.}
\end{qabox}

\begin{qabox}[Hochreiter]{How can you do feature selection?}
\ans Three families. Filter: rank features by a score like correlation or mutual information.
Wrapper: search subsets with the model, e.g.\ forward or backward selection. Embedded: the model
selects while it trains, e.g.\ L1 regularization (Lasso) pushes weights to exactly zero.
\end{qabox}

\begin{qabox}[Klambauer]{Why initialize randomly? What does good initialization do?}
\ans Equal weights make all units in a layer learn the same thing, so we break the symmetry with
random weights. The scale should keep the variance of activations forward and of deltas backward
constant, so nothing vanishes or explodes.
\formula $\mathrm{Var}(w)=1/n_{\mathrm{in}}$ (LeCun), $2/n_{\mathrm{in}}$ (He, ReLU),
$2/(n_{\mathrm{in}}+n_{\mathrm{out}})$ (Glorot).
\src{DLNN script WS24, Ch.~9, PDF p.~123--129.}
\end{qabox}

\begin{qabox}[Lehner]{Universal approximation theorem.}
\ans One hidden layer with finitely many units and a non-polynomial activation can approximate
any continuous function on a compact set arbitrarily well. It says such a network exists; it does
not say how many units, nor that gradient descent finds it.
\src{DLNN I+II script SS25, ``Deep Learning theory'', PDF p.~321.}
\end{qabox}

\begin{qabox}[Klambauer]{CNNs: why do they work for images? What is ResNet?}
\ans Local filters with shared weights slide over the image: few parameters, translation
equivariance. Pooling reduces size and adds invariance. Depth grows the receptive field. ResNet
adds identity skips, $\mathbf{a}^{[l+1]}=\mathbf{a}^{[l]}+F(\mathbf{a}^{[l]})$, so the Jacobian is
$I+\partial F$ and gradients pass through hundreds of layers.
\formula Output size $\lfloor(n+2p-k)/s\rfloor+1$.
\src{DLNN script WS24, Ch.~6; SS25 script PDF p.~176, 323.}
\end{qabox}
